% =====================================================================
%  chunks/ch04i.tex —— 原书 p.239–243（PDF 246–250）
%  第 4 章 §4（下）：整 Pontryagin 类、法映射与法配边、S-对偶、光滑与 PL 结构
% =====================================================================
% [原文 p.239]
Cartan--Serre 理论表明，稳定群 $\pi_{n+N}(T\xi)$ 具有形式
\begin{equation}
\pi_{n+N}(T\xi)=\Z\oplus\bar{\mathcal{N}}(X^n,\xi) \tag{4.14}
\end{equation}
其中 $\bar{\mathcal{N}}(X^n,\xi)$ 是有限阿贝尔群；关于这个分解，上述度数为 $1$ 的法映射的法配边类表示如下：
\begin{equation}
\alpha\in \mathcal{N}(X^n,\xi),\qquad \alpha=1+a,\qquad a\in\bar{\mathcal{N}}(X^n,\xi). \tag{4.15}
\end{equation}

在 $X^n=M^n$ 是光滑流形、$\xi$ 是其稳定法丛 $\nu_M$（$M^n\subset\R^{n+N}$）的情形，$\pi_{n+N}(T\xi)$ 的单位元 $1$ 对应于恒等映射对
\[
\begin{tikzpicture}[>=Stealth,baseline=(current bounding box.center)]
\node (a) at (0,1.15) {$\nu_M$};
\node (b) at (2.2,1.15) {$\nu_M$};
\node (c) at (0,0) {$M^n$};
\node (d) at (2.2,0) {$M^n$};
\draw[->] (a)--(b) node[midway,above] {$\id$};
\draw[->] (a)--(c);
\draw[->] (b)--(d);
\draw[->] (c)--(d) node[midway,below] {$\id$};
\end{tikzpicture}
\]
并且下列断言成立：

\begin{enumerate}[label=\arabic*.,leftmargin=2.2em,itemsep=4pt]
\item 设 $f_1:M^n_1\to M^n$、$f_2:M^n_2\to M^n$ 是度数为 $1$ 的映射，它们连同法丛的诱导映射给出同一个法配边类。若映射 $f_1$、$f_2$ 还都是同伦等价，则存在同伦 $n$-球面 $\Sigma^n\in bP^{n+1}$，使
\begin{equation}
M^n_1\cong\Sigma^n\#M^n_2. \tag{4.16}
\end{equation}
由此可知：$n\ge5$ 时流形 $M^n_1$ 与 $M^n_2$ 同胚，并且从两者中各去掉一点所得的带孔流形实际上微分同胚；而当 $n$ 为大于 $5$ 的偶数时，$M^n_1$ 与 $M^n_2$ 实际上微分同胚（Novikov）。

\item 下面的结果适用于 $M^n$ 上任一给定的向量丛 $\xi$：$n\ge5$ 为奇数时，每个法配边类 $\alpha\in\mathcal{N}(M^n,\xi)$ 都可由某个同伦等价 $f:M'\to M^n$ 实现。向量丛 $\xi$ 的 Pontryagin 类满足如下条件：
\[
L(p_1(\xi),\dots,p_k(\xi))=\tau(M^n)
\]
这一条件对 $n=4k$ 的情形也是充分必要的（Browder，Novikov）。
\end{enumerate}
% [Tutor Note: 式 (4.15) 里的 $1$ 是同伦群 $\pi_{n+N}(T\xi)$ 的单位元，不是整数；$\mathcal{N}(X^n,\xi)$ 的加法结构由 Thom 空间的 $H$-空间乘法给出，几何上对应法配边的连通和。把 $\alpha$ 拆成 $1+a$ 正是手术理论的出发点：整分量由式 (4.13) 的 Hurewicz 同构固定，全部几何信息都留在有限群 $\bar{\mathcal{N}}$ 里。]

% [原文 p.240]
由第一个结果立即得到一个推论：对每个给定的同伦型，只有有限多个两两不微分同胚的闭单连通 $n$ 维流形（$n\ge5$）具有相同的 Pontryagin 类 $p_j\in H^{4j}(M^n;\Q)$。利用 Pontryagin 类的拓扑不变性，由此又可推出：每个闭单连通的 $n$ 维拓扑流形（$n\ge5$）上只存在有限多个互不相同的光滑结构。由于这一理论对 $PL$ 流形同样有效，还可知任何这样的拓扑流形上可定义的 $PL$ 结构也只有有限多个。
% [Tutor Note: 这里只用到 $p_j(M^n)$ 的拓扑不变性，不需要任何微分结构；因此「同一拓扑流形上光滑结构有限」与「同一同伦型内流形有限」是同一套配边计算的两种读法。该推论要求 $n\ge5$，因为证明中的 Morse 手术要用到 $n\ge5$ 时的 Whitney 技巧。]

\begin{remark}{}
这里顺便注意：对 $PL$ 流形，第一个结果的类似物给出的结论有偶然的加强，即流形 $M^n_1$ 与 $M^n_2$ 实际上将是 $PL$ 同胚的——原因在于所有同伦 $n$-球面都 $PL$ 同胚于 $S^n$。\hfill$\square$
\end{remark}

对满足上面「结果 1」各项假设的流形（即同属一个同伦类，且到 $M^n$ 的映射延拓到稳定法丛后给出同一个法配边类），其微分同胚类的精确分类这样得到：取集合 $\mathcal{N}(M^n,\nu_n)$ 中形如 $1+a\in\pi_{N+n}(T\nu_n)$ 的元素，再模掉子群 $\mathcal{G}(M^n,\nu_n)$——后者由代表度数为 $1$ 的自同构（即保稳定法丛 $f^*\nu_n=\nu_n$ 的度数为 $1$ 的自映射 $f:M^n\to M^n$）的同伦类的那些元素组成。于是，满足「结果 1」各项假设的流形的微分同胚类，与
\begin{equation}
\mathcal{N}(M^n,\nu_n)/\mathcal{G}(M^n,\nu_n) \tag{4.17}
\end{equation}
的陪集自然一一对应。

还有一个小一些的群 $\mathcal{J}(M^n,\nu_n)$，它是自同构群 $\mathcal{G}(M^n,\nu_n)$ 的子群，由丛 $\nu_n$ 的、在底流形 $M^n$ 上诱导恒等映射的那些自同构组成。轨道集
\begin{equation}
\mathcal{N}(M^n,\nu_n)/\mathcal{J}(M^n,\nu_n) \tag{4.18}
\end{equation}
与下述问题有关：给定度数为 $1$ 的法同伦等价 $f:M^n_1\to M^n$，何时 $f$ 可形变为微分同胚（更确切地说，可形变到「与 $bP^{n+1}$ 中某个 Milnor 球面 $\Sigma^n$ 作连通和」的意义下的微分同胚）？答案是：这样的映射 $f$ 所定义的法配边可如此形变，当且仅当它在陪集 (4.18) 上的投影，与那个特殊元素（由恒等映射 $M^n\to M^n$ 代表）的投影相同。此外，若 $f_1:M^n_1\to M^n$ 与 $f_2:M^n_2\to M^n$ 是度数为 $1$ 的法映射（代表法配边类），则它们对应 (4.18) 中同一陪集，当且仅当 $f_1f_2^{-1}$ 同伦于一个微分同胚（模 $bP^{n+1}$）。

% [原文 p.241]
在给定向量丛 $\xi$ 是 $M^n\subset S^{N+n}$ 的稳定法丛、即 $\xi=\nu_n$ 的情形，Thom 空间 $T\nu_n$ 有特别的几何意义：取适当小的 $\eps$，令 $U_\eps$ 是流形 $M^n\subset S^{N+n}$ 的 $\eps$-邻域；则按定义
\[
T\nu_M=S^{N+n}\big/(S^{N+n}\setminus U_\eps),
\]
即把邻域 $U_\eps$ 的补集等同为一点后所得的 $S^{N+n}$ 的商空间（这里 $N\gg n$）。相应的自然映射 $S^{N+n}\to T\nu_n$ 在 $U_\eps$ 上限制为恒等映射；它的度数为 $1$，并且（在 $\pi_{n+N}(T\nu_n)$ 中）代表流形 $M^n$。当 $n$ 具有 $n=4k+2$ 的形式时，$\pi_{n+N}(T\nu_n)$ 中形如 $1+a$ 的元素（其中 $a$ 取遍 $\bar{\mathcal{N}}(M^n,\nu_n)$ 中某个指数为 $1$ 或 $2$ 的子群 $B$）恰好就是那些对应于与 $M^n$ 同伦型的 $n$ 维流形的元素。（空间 $T\nu_n$ 常常便于对具体流形作计算。）

上述理论已被 Golo 与 Wall（20 世纪 60 年代中期）推广到带边流形的情形。

20 世纪 60 年代中期，好几位作者\footnote{A. Schwarz，Novikov，Sullivan。}展示了「$S$-对偶」概念的作用，这一概念带有更一般的范畴性质。Spanier 与 G. W. Whitehead 把嵌入充分高维球面 $S^m$ 中的 $CW$-复形 $K$ 的 $S$-对偶 $DK$ 就定义为补集 $S^m\setminus K$；于是算子 $D$ 作用在 $CW$-复形 $K$ 上如下：
\[
D:K\mapsto S^m\setminus K.
\]
特别地，球面的 $S$-对偶（同伦等价意义下）是球面。在映射 $f:K\to L$（的同伦类）上，$S$-对偶算子把箭头反转（$f$ 是嵌入时这一点很清楚；一般情形借映射柱化归为此）：
\[
Df:DL\longrightarrow DK,
\]
\[
[K,L]_s=[DL,DK]_s,
\]
其中 $[\,,]_s$ 表示映射的稳定同伦类集合。20 世纪 60 年代初 Atiyah 注意到：闭流形 $M^n$ 的稳定法丛的 Thom 空间的 $S$-对偶，是 $M^n$ 上平凡向量丛 $\eps^q$ 的 Thom 空间：
\[
D(T\nu_M)=\Sigma(M^n_+)=T\eps^q,
\]
其中 $M^n_+=M^n\cup\{*\}$（与一点空间的不交并）。映射 $f:S^{N+n}\to T\nu_M$ 的 $S$-对偶原来是映射 $Df:T\eps^q\to S^q$，即从平凡向量丛 $\eps^q$ 的 Thom 空间到 $S^q$、在每条纤维上度数为 $1$ 且随底点参数变化的映射。
% [Tutor Note: $S$-对偶把 $T\nu_M$ 换成 $T\eps^q$，是整套理论的枢纽：它把「法映射」这一几何对象翻译成从 $M^n$ 到 $SG$ 的映射，于是法配边类成了同伦类，分类问题化为计算 $\pi_j(SG/SO)$。下面的商 (4.17)、(4.18) 正是这一翻译的产物。]

% [原文 p.242]
用 $SG_q$ 记度数为 $1$ 的映射 $S^q\to S^q$ 所成的半群。对应于每个法配边类 $\alpha\in\mathcal{N}(M^n,\xi)$，有一个从底流形 $M^n$ 到半群 $SG_q$ 的映射：
\[
\bar Df:M^n\longrightarrow SG_q,\qquad (q\to\infty).
\]
此映射在商 $SG_q/SO$ 中的像，便构成类 $\alpha$ 中各个流形 $M^n_1$ 之间存在微分同胚的障碍。注意 $q\to\infty$ 时，$SG_q$ 的同伦性质（计及其 $H$-空间结构）会稳定下来。自然嵌入给出谱
\[
SG_q\subset SG_{q+1}\subset\cdots
\]
其极限记为 $SG$。还有嵌入
\[
SG\supset SPL\supset SO,
\]
其中 $SPL$ 表示充分高维圆盘上固定原点、保持定向的 $PL$-自同构的芽所成的群。同伦群的诱导单同态 $\pi_j(SO)\to\pi_j(SG)$ 与 Whitehead 同态 $J$（见上文 §3）一致；这里用到了同构
\[
\pi_j(SG)=\pi_{N+j}(S^N),\qquad N>j+1
\]
所作的等同。

球丛（即以球面为纤维的丛）的万有底空间 $BSG$ 这样赋予意义：用与投影交换、且在纤维上度数为 $1$ 的丛映射来定义这类丛的等价。空间 $SG_q$ 显然具有 $\Omega^q(S^q)$（$q$ 维环路空间的单位元连通分支）的同伦型，因此
\[
\pi_j(\Omega^q(S^q))\cong\pi_{q+j}(S^q)\cong\pi_j(SG),
\]
但这两个空间的乘法结构并不相同。注意空间 $BSG$ 不具有任何 $(q-1)$-重环路空间 $\Omega^{q-1}(Y)$ 的同伦型。

这套理论几乎原样搬到 $PL$ 流形的情形。流形的切丛与稳定法丛，配以维数不定的圆盘上固定原点的 $PL$-自同构的芽群，按 Milnor 的叫法称为\emph{微丛}（microbundle）。20 世纪 60 年代初由好几位作者证明的一个定理断言：$PL$ 流形具有相容的光滑结构，当且仅当其稳定法（或切）微丛可约化为向量 $O$-丛，$O\subset PL$（可定向情形下可约化为向量 $SO$-丛，$SO\subset SPL$）。表示 $PL$ 微丛可约化为光滑微丛之障碍的相对同伦群 $\pi_j(SPL,SO)$，同构于相应的群 $\Gamma^{j-1}$：
\begin{equation}
\pi_j(SPL,SO)=\pi_j(SPL/SO)=\Gamma^{j-1}. \tag{4.19}
\end{equation}

% [原文 p.243]
由此可得：$PL$ 流形 $M$ 的 $j$ 维骨架的邻域上存在光滑结构的障碍，是上同调类 $H^j(M;\Gamma^j)$。$PL$ 流形的理论从 20 世纪 30 年代起由许多作者（Newman、Whitehead、Zeeman、Munkres、Mazur、Hirsch 等人）发展起来。

在 $PL$ 流形上切丛平凡的小开区域，总可以定义「局部的」光滑结构；正是这一事实使得上述关于给定同伦型光滑流形分类理论的整套技术手段能自然地搬到 $PL$ 流形的情形：因为消去相应同伦群的核的 Morse 手术（见上文）都是在小的可平行化区域上进行的。

不过还需计及一个事实：一般而言，$PL$-向量丛在底空间上诱导恒等映射的自同构比光滑（或 $SO$-）丛要多。因此，在 $S$-对偶的途径中，流形之间存在 $PL$-同构的障碍由映射
\[
M^n\longrightarrow SG/SPL
\]
的同伦类代表，而非由到 $SG/SO$ 的映射代表。1966 年 Sullivan 与 Wagoner 证明：群 $\pi_j(SG/SPL)$（或与之同构的群 $\pi_j(G/PL)$）结构简单，并服从漂亮的 4-周期性：
\[
\pi_j(SG/SPL)=\pi_j(G/PL)=
\begin{cases}
\Z, & j=4k,\\
0, & j=4k+1,\\
\Z/2, & j=4k+2,\\
0, & j=4k+3.
\end{cases}
\]
随后，Sullivan 利用空间 $SG/SPL$ 的同伦结构\footnote{空间 $G/PL$、$SG/SPL$ 的同伦性质，在 Madsen、Milgram 的著作 \emph{Classifying Spaces in Surgery and Cobordism of Manifolds}（Princeton University Press, 1979, Annals of mathematics studies, no. 92）中有完整描述。}及其与 Bott 周期性的联系，连同其他性质，确立了一些重要的一般事实；然而直到现在（1994 年），Sullivan 结果的完整证明尚未公开发表。

如前所述，同伦群 $\pi_j(SPL/SO)$ 同构于群 $\Gamma^{j-1}$。上述结果表明，包含同态
\[
\pi_j(SO)\longrightarrow\pi_j(SPL)
\]
在 $j\le6$ 时实际上是同构。当 $j=7$ 时，$\pi_7(SO)$ 在 $\pi_7(SPL)$ 中的像的阶可被 $7$ 整除。后一结果不难建立：转到万有 $SO$-丛与 $SPL$-丛的底空间之间的包含 $BSO\to BSPL$，并利用 $PL$ 流形的 Pontryagin 类可以是分母为 $7$ 的分数（见 §3）。
% [Tutor Note: 式 (4.19) 记 $\pi_j(SPL/SO)=\Gamma^{j-1}$，而本段又说 $j$ 维骨架上的障碍属于 $H^j(M;\Gamma^j)$，两处 $\Gamma$ 的指标相差 $1$。这是原书对 $\Gamma$ 群编号的疏失，计算时以式 (4.19) 的显式同构为准，不要按 $H^j$ 里的上标去反推 (4.19)。]

由此可推出如下推论：

% CHUNK ch04i DONE
